% Source: https://artofproblemsolving.com/community/c3h2501621_2015_mathcounts_state_problem_28
% Stable identity: aops-topic-2501621-q01
Scale the pentagon to side $1$ and let $F$ be the intersection of the diagonals $AC$ and $BE$.

\begin{asy}
usepackage("amsmath");
size(5cm);
defaultpen(fontsize(10pt));
pair A=dir(90), B=dir(162), C=dir(234), D=dir(306), Ep=dir(18);
pair M=(A+Ep)/2;
pair Fq=extension(A,C,B,Ep), Z=extension(A,C,B,M);
draw(A--B--C--D--Ep--cycle,linewidth(0.9));
draw(A--C,blue); draw(B--Ep,gray); draw(B--M,red);
dot(M); dot(Fq); dot(Z);
label("$A$",A,N); label("$B$",B,W); label("$C$",C,SW); label("$D$",D,SE); label("$E$",Ep,E);
label("$M$",M,NE); label("$F$",Fq,SE); label("$Z$",Z,NW);
\end{asy}

\textbf{Lengths in the pentagon.} Each interior angle is $108^\circ$, and the isosceles triangles $ABC$ and $EAB$ have base angles $36^\circ$, so $\angle BAC=\angle ABE=36^\circ$. In triangle $ABF$ these two angles are equal, so $AF=BF$, and $\angle AFB=108^\circ$. In triangle $BCF$, $\angle CBF=108^\circ-36^\circ=72^\circ$ and $\angle BFC=180^\circ-108^\circ=72^\circ$, so $CF=CB=1$. Let $d=AC=BE$ be the diagonal; then $AF=BF=d-1$ and $FE=BE-BF=1$.

Triangles $ABF$ and $ACB$ have the same angles ($36^\circ,36^\circ,108^\circ$), so $\dfrac{AF}{AB}=\dfrac{AB}{AC}$, i.e. $d-1=\dfrac1d$. Thus $d^2-d-1=0$ and $d=\varphi=\dfrac{1+\sqrt5}{2}$, with $AF=BF=\varphi-1=\dfrac1\varphi$ and $FE=1$.

\textbf{Mass points in triangle $ABE$.} In triangle $ABE$, the segment $BM$ is a median and $AF$ is a cevian with $\dfrac{BF}{FE}=\dfrac{1}{\varphi}$; they meet at $Z$. Mass points place weights at the vertices so that each point on a side is the balance point (centre of mass) of the weights at the ends of that side, like a see-saw: weights $w_1$, $w_2$ at the ends balance at a point that divides the segment in the ratio $w_2:w_1$. Then the intersection of two cevians is the balance point of all three weights, and it divides each cevian inversely to the weights at its ends.
\begin{itemize}
\item $M$ is the midpoint of $AE$, so put weight $1$ at $A$ and $1$ at $E$.
\item $F$ balances $B$ and $E$ with $\dfrac{BF}{FE}=\dfrac1\varphi$, so the weight at $B$ is $\varphi$ times the weight at $E$: weight $\varphi$ at $B$, and $F$ carries $\varphi+1=\varphi^2$.
\item On the cevian $AF$, $Z$ balances weight $1$ at $A$ against $\varphi^2$ at $F$, so $\dfrac{AZ}{ZF}=\varphi^2$.
\end{itemize}
(Menelaus' theorem in triangle $AFE$ with the line $B$--$Z$--$M$ confirms this: $\dfrac{AZ}{ZF}\cdot\dfrac{FB}{BE}\cdot\dfrac{EM}{MA}=1$ gives $\dfrac{AZ}{ZF}=\dfrac{BE}{FB}=\dfrac{\varphi}{\varphi-1}=\varphi^2$.)

\textbf{Finish.} Hence
\[AZ=\dfrac{\varphi^2}{\varphi^2+1}\,AF=\dfrac{\varphi^2}{\varphi^2+1}\cdot\dfrac1\varphi=\dfrac{\varphi}{\varphi^2+1}=\dfrac{\varphi}{\varphi+2}=\dfrac{1+\sqrt5}{5+\sqrt5}=\dfrac{(1+\sqrt5)(5-\sqrt5)}{20}=\dfrac{\sqrt5}{5}.\]
So in any regular pentagon $ZA=\dfrac{AB}{\sqrt5}$, and $ZA=3$ gives $AB=\boxed{3\sqrt5}$.
